
\documentclass[12pt,a4paper]{article}

\usepackage[a4paper,margin=1in]{geometry}
\usepackage{graphicx}
\usepackage{booktabs}
\usepackage{amsmath,amssymb}
\usepackage{float}
\usepackage{hyperref}
\usepackage{xcolor}
\usepackage{listings}
\usepackage{enumitem}
\usepackage{fancyhdr}
\usepackage{longtable}

\hypersetup{colorlinks=true,linkcolor=blue,urlcolor=blue}

\pagestyle{fancy}
\fancyhf{}
\lhead{ICS1512}
\rhead{Machine Learning Algorithms Laboratory}
\cfoot{\thepage}

\lstset{
language=Python,
basicstyle=\ttfamily\small,
numbers=left,
frame=single,
breaklines=true,
showstringspaces=false
}

\begin{document}

\begin{tabular}{|p{4cm}|p{11cm}|}
\hline
Experiment No. & \\ \hline
Experiment Title & \footnotemark\\ \hline
Student Name & \\ \hline
Register Number & \\ \hline
Faculty & Dr. Poreddy Ajay Kumar Reddy \\ \hline
Submission Date & \\ \hline
\end{tabular}
\footnotetext{\textbf{GitHub Repository: }\url{https://github.com/username/repository}}

\section{Objective}
State the objective of the experiment.

\section{Problem Statement}
Describe the problem, input, output, and objective.

\section{Dataset Description}
\begin{longtable}{|p{6cm}|p{8cm}|}
\hline
Dataset Name & \\ \hline
Dataset Source & \\ \hline
Number of Samples & \\ \hline
Number of Features & \\ \hline
Number of Classes & \\ \hline
Missing Values & \\ \hline
Train-Test Split & \\ \hline
\end{longtable}

\section{Exploratory Data Analysis}
Include all relevant visualizations.
\begin{itemize}
\item Dataset overview
\item Statistical summary
\item Missing value analysis
\item Class distribution
\item Correlation matrix
\item Feature distribution
\end{itemize}

\subsection*{Figure Template}

\begin{figure}[H]
\centering
\includegraphics[width=0.7\linewidth]{example-image}
\caption{Replace with your figure caption}
\end{figure}

\textbf{Inference:}

Write 3--5 sentences explaining what the graph indicates and how it helps in understanding the dataset.

\section{Data Preprocessing}

Explain:
\begin{itemize}
\item Missing value handling
\item Encoding
\item Feature scaling
\item Train-test split
\item Feature engineering (if applicable)
\end{itemize}

\section{Machine Learning Algorithm}
Explain:
\begin{itemize}
\item Working principle
\item Mathematical formulation
\item Advantages
\item Limitations
\end{itemize}

\section{Experimental Setup}
\begin{tabular}{ll}
\toprule
Python Version & \\
Scikit-Learn Version & \\
IDE & \\
Operating System & \\
Hardware & \\
\bottomrule
\end{tabular}

\section{Hyperparameter Selection}

\begin{tabular}{ll}
\toprule
Parameter & Value\\
\midrule
 & \\
 & \\
\bottomrule
\end{tabular}



\section{Results}

\begin{tabular}{lc}
\toprule
Metric & Value\\
\midrule
Accuracy & \\
Precision & \\
Recall & \\
F1-score & \\
ROC-AUC & \\
\bottomrule
\end{tabular}

\section{Result Visualization}

Include all graphs generated during the experiment.

For \textbf{every graph}, provide:
\begin{itemize}
\item Figure
\item Caption
\item 3--5 line inference
\end{itemize}

Recommended figures:
\begin{enumerate}
\item Correlation Matrix
\item Class Distribution
\item Confusion Matrix
\item ROC Curve
\item Precision-Recall Curve
\item Learning Curve (if applicable)
\item Feature Importance (if applicable)
\end{enumerate}

\section{Discussion}
Discuss observations, strengths, limitations and improvements.

\section{Conclusion}
Summarize the experiment and key findings.

\section{References}
Use IEEE style.

\end{document}
